\chapter*{ABSTRACT}
\addcontentsline{toc}{chapter}{ABSTRACT}

The rapid growth of Vietnamese online news has created two parallel needs:
fast summarization for readers who skim and reliable fact-checking for readers who
verify. Summarizers built on a single Large Language Model (LLM)
still produce uneven outputs, and Mixture-of-Agents (MoA) --- a recent technique
that combines several LLMs through a proposer--aggregator pipeline ---
has not yet been studied for Vietnamese news summarization. At the same time, the metrics most often
used to evaluate summarization in Vietnamese research (ROUGE, BLEU, BERTScore)
measure $n$-gram or embedding overlap against the source article, which is a poor
indicator of summary quality when the system is allowed to rewrite and paraphrase.

This thesis makes 3 contributions. First, it presents \textbf{Fiber}, a
browser extension and Next.js backend that summarize and fact-check Vietnamese
news in real time, supporting LLMs across OpenAI, Anthropic, Gemini, and even HuggingFace. Second, it adapts the Mixture-of-Agents pipeline of Wang et al.\ (2024)
to grounded Vietnamese summarization, introducing a source-article residual
connection into the aggregator prompt so that synthesis remains faithful to the
input article. Third, and most importantly, it proposes a
\textbf{three-axis evaluation framework} --- content retention (Axis~A: ROUGE,
BLEU, BERTScore), automatic quality and preference (Axis~B: rubric, pairwise
LLM-judge, factuality entailment), and human validation (Axis~C: blind $K$-way
ranking with inter-rater agreement) --- and uses it to show why
overlap metrics alone cannot evaluate summarization.

On a topic-balanced dataset of 154 Vietnamese news articles from
\textit{tienphong.vn}, the MoA pipeline beats the strongest single-model
baseline (best draft selection) in $93$ of $135$ decisive pairwise verdicts
($68.9\%$, sign-test $p < 0.0001$); it consistently outperforms individual proposers. Rubric scoring and overlap
metrics both agree with this verdict. A comprehensive human study of 17 raters across
all $48$ tasks (n=192 rankings) reveals a $28$-percentage-point gap between the
LLM judge and human preference: humans pick the fused output $40.9\%$ of the
time, not $68.9\%$. This gap, together with a topic-conditioned interaction
(humans favor short proposer drafts on bureaucratic content but the fused
summary on entertainment content), is the methodological finding of the thesis.

\vspace{0.4em}
\textbf{Keywords:} Vietnamese summarization, fact-checking, large language
models, Mixture-of-Agents, LLM-as-judge, browser extension, three-axis
evaluation.
